\section{转轴和转角}

\subsection{用转轴和转角表示的空间运动刚体角速度}

用转轴$\boldsymbol{n}$（单位矢量）和转角$\theta$表示的方向余弦矩阵为
\begin{equation}
A_{ij} = I_3\cos\theta + \tilde{\boldsymbol{n}}\sin\theta + \boldsymbol{n}\boldsymbol{n}^{\mathrm{T}}(1 - \cos\theta)
\label{eq:7.1}
\end{equation}

用转轴和转角表示的相对转动角速度矢量在动系（$j$系）中的投影列阵为
\begin{equation}
\boldsymbol{\omega}_{ij|j} = \boldsymbol{n}\dot{\theta} + \sin\theta\,\dot{\boldsymbol{n}} - (1 - \cos\theta)\,\tilde{\boldsymbol{n}}\dot{\boldsymbol{n}}
\label{eq:7.2}
\end{equation}

上式两端同时左乘$\boldsymbol{n}^{\mathrm{T}}$可得
\begin{equation}
\dot{\theta} = \boldsymbol{n}^{\mathrm{T}}\boldsymbol{\omega}_{ij|j}
\end{equation}

回代后可得
\begin{equation}
\dot{\boldsymbol{n}} = \frac{1}{2\sin\theta}\bigl[\sin\theta\,I_3 - (1 + \cos\theta)\,\tilde{\boldsymbol{n}}\bigr]\,\tilde{\boldsymbol{n}}\,\boldsymbol{\omega}_{ij|j}
= \frac{1}{2}\bigl[I_3 - \mathrm{ctan}(\theta/2)\,\tilde{\boldsymbol{n}}\bigr]\,\tilde{\boldsymbol{n}}\,\boldsymbol{\omega}_{ij|j}
\label{eq:7.3}
\end{equation}

\section{欧拉参数}

方向余弦矩阵与转轴和转角的关系式为
\begin{equation}
A_{ij} = I_3\cos\theta + \tilde{\boldsymbol{n}}\sin\theta + \boldsymbol{n}\boldsymbol{n}^{\mathrm{T}}(1 - \cos\theta)
\end{equation}

利用三角函数的倍角公式：
\begin{equation}
\cos\theta = 2\cos^2\frac{\theta}{2} - 1, \quad
\sin\theta = 2\cos\frac{\theta}{2}\sin\frac{\theta}{2}, \quad
1 - \cos\theta = 2\sin^2\frac{\theta}{2}
\end{equation}

定义四个欧拉参数：
\begin{equation}
\begin{bmatrix}
\varepsilon_1 & \varepsilon_2 & \varepsilon_3
\end{bmatrix}^{\mathrm{T}}
\triangleq \boldsymbol{n}\sin\frac{\theta}{2} =: \bar{\boldsymbol{\varepsilon}}, \quad
\varepsilon_4 \triangleq \cos\frac{\theta}{2}
\label{eq:7.4}
\end{equation}

则有如下多项式矩阵形式的方向余弦矩阵：
\begin{equation}
A_{ij} = I_3(2\varepsilon_4^2 - 1) + 2\tilde{\bar{\boldsymbol{\varepsilon}}}\,\varepsilon_4 + 2\bar{\boldsymbol{\varepsilon}}\bar{\boldsymbol{\varepsilon}}^{\mathrm{T}}
\end{equation}

欧拉参数满足如下约束方程：
\begin{equation}
\bar{\boldsymbol{\varepsilon}}^2 + \varepsilon_4^2
= n^2\sin^2\frac{\theta}{2} + \cos^2\frac{\theta}{2}
= \sin^2\frac{\theta}{2} + \cos^2\frac{\theta}{2} = 1
\end{equation}

用欧拉参数表示的方向余弦矩阵形式不唯一：
\begin{equation}
\begin{cases}
A_{ij} = I_3(2\varepsilon_4^2 - 1) + 2\tilde{\bar{\boldsymbol{\varepsilon}}}\,\varepsilon_4 + 2\bar{\boldsymbol{\varepsilon}}\bar{\boldsymbol{\varepsilon}}^{\mathrm{T}}, \\[4pt]
A_{ij} = I_3(1 - 2\bar{\boldsymbol{\varepsilon}}^2) + 2\tilde{\bar{\boldsymbol{\varepsilon}}}\,\varepsilon_4 + 2\bar{\boldsymbol{\varepsilon}}\bar{\boldsymbol{\varepsilon}}^{\mathrm{T}}
\end{cases}
\label{eq:7.5}
\end{equation}

用欧拉参数表示的方向余弦矩阵的展开形式为
\begin{equation}
A_{ij} = I_3(1 - 2\bar{\boldsymbol{\varepsilon}}^2) + 2\tilde{\bar{\boldsymbol{\varepsilon}}}\,\varepsilon_4 + 2\bar{\boldsymbol{\varepsilon}}\bar{\boldsymbol{\varepsilon}}^{\mathrm{T}}
=
\begin{bmatrix}
1 - 2\varepsilon_2^2 - 2\varepsilon_3^2 & 2(\varepsilon_1\varepsilon_2 - \varepsilon_3\varepsilon_4) & 2(\varepsilon_1\varepsilon_3 + \varepsilon_2\varepsilon_4) \\[4pt]
2(\varepsilon_2\varepsilon_1 + \varepsilon_3\varepsilon_4) & 1 - 2\varepsilon_3^2 - 2\varepsilon_1^2 & 2(\varepsilon_2\varepsilon_3 - \varepsilon_1\varepsilon_4) \\[4pt]
2(\varepsilon_3\varepsilon_1 - \varepsilon_2\varepsilon_4) & 2(\varepsilon_3\varepsilon_2 + \varepsilon_1\varepsilon_4) & 1 - 2\varepsilon_1^2 - 2\varepsilon_2^2
\end{bmatrix}
\label{eq:7.6}
\end{equation}

四个欧拉参数满足模为1的约束方程，即
\begin{equation}
\varepsilon_1^2 + \varepsilon_2^2 + \varepsilon_3^2 + \varepsilon_4^2 = 1
\label{eq:7.7}
\end{equation}

由$\tilde{\boldsymbol{\omega}}_{ij|j} = A_{ij}^{\mathrm{T}}\dot{A}_{ij}$可得用欧拉参数及其导数表示的$\boldsymbol{\omega}_{ij|j}$的表达式。

用欧拉参数及其导数表示的$\boldsymbol{\omega}_{ij|j}$的表达式为
\begin{equation}
\boldsymbol{\omega}_{ij|j} = 2
\begin{bmatrix}
\varepsilon_4  & \varepsilon_3  & -\varepsilon_2 & -\varepsilon_1 \\[2pt]
-\varepsilon_3 & \varepsilon_4  & \varepsilon_1  & -\varepsilon_2 \\[2pt]
\varepsilon_2  & -\varepsilon_1 & \varepsilon_4  & -\varepsilon_3
\end{bmatrix}
\begin{bmatrix}
\dot{\varepsilon}_1 \\ \dot{\varepsilon}_2 \\ \dot{\varepsilon}_3 \\ \dot{\varepsilon}_4
\end{bmatrix}
= 2
\begin{bmatrix}
\varepsilon_4 I_3 - \tilde{\bar{\boldsymbol{\varepsilon}}} & -\bar{\boldsymbol{\varepsilon}}
\end{bmatrix}
\dot{\boldsymbol{\varepsilon}} =: H_{\mathrm{EP}}\,\dot{\boldsymbol{\varepsilon}}
\label{eq:7.8}
\end{equation}

进一步可得$\boldsymbol{\omega}_{ij|j}$在$i$系中的投影：
\begin{equation}
\boldsymbol{\omega}_{ij|i} = A_{ij}\boldsymbol{\omega}_{ij|j} = 2
\begin{bmatrix}
\varepsilon_4 I_3 + \tilde{\bar{\boldsymbol{\varepsilon}}} & -\bar{\boldsymbol{\varepsilon}}
\end{bmatrix}
\dot{\boldsymbol{\varepsilon}}
\end{equation}

以及
\begin{equation}
\frac{\partial\boldsymbol{\omega}_{ij|i}}{\partial\boldsymbol{\varepsilon}} = 2
\begin{bmatrix}
-\dot{\varepsilon}_4  & -\dot{\varepsilon}_3  & \dot{\varepsilon}_2  & \dot{\varepsilon}_1  \\[2pt]
\dot{\varepsilon}_3   & -\dot{\varepsilon}_4  & -\dot{\varepsilon}_1 & \dot{\varepsilon}_2  \\[2pt]
-\dot{\varepsilon}_2  & \dot{\varepsilon}_1   & -\dot{\varepsilon}_4 & \dot{\varepsilon}_3
\end{bmatrix}
= 2
\begin{bmatrix}
-\dot{\varepsilon}_4 I_3 + \tilde{\dot{\bar{\boldsymbol{\varepsilon}}}} & \dot{\bar{\boldsymbol{\varepsilon}}}
\end{bmatrix}
\label{eq:7.9}
\end{equation}

\subsection*{推导过程}

取$\boldsymbol{\omega}_{ij|j}$的$x$分量：
\begin{align}
(\omega_x)_{ij|j} &= a_{13}\dot{a}_{12} + a_{23}\dot{a}_{22} + a_{33}\dot{a}_{32} \nonumber\\
&= \bigl[4\varepsilon_1(\varepsilon_1\varepsilon_4 - \varepsilon_2\varepsilon_3) + 2\varepsilon_4\bigr]\dot{\varepsilon}_1
   + \bigl[4\varepsilon_2(\varepsilon_1\varepsilon_4 - \varepsilon_2\varepsilon_3) + 2\varepsilon_3\bigr]\dot{\varepsilon}_2 \nonumber\\
&\quad + \bigl[4\varepsilon_3(\varepsilon_1\varepsilon_4 - \varepsilon_2\varepsilon_3) - 2\varepsilon_2\bigr]\dot{\varepsilon}_3
   + \bigl[4\varepsilon_4(\varepsilon_1\varepsilon_4 - \varepsilon_2\varepsilon_3) - 2\varepsilon_1\bigr]\dot{\varepsilon}_4 \nonumber\\
&= 2\varepsilon_4\dot{\varepsilon}_1 + 2\varepsilon_3\dot{\varepsilon}_2 - 2\varepsilon_2\dot{\varepsilon}_3 - 2\varepsilon_1\dot{\varepsilon}_4
\end{align}

其中$A_{ij}$展开为
\begin{equation}
A_{ij} =
\begin{bmatrix}
1 - 2\varepsilon_2^2 - 2\varepsilon_3^2 & 2(\varepsilon_1\varepsilon_2 - \varepsilon_3\varepsilon_4) & 2(\varepsilon_1\varepsilon_3 + \varepsilon_2\varepsilon_4) \\[4pt]
2(\varepsilon_2\varepsilon_1 + \varepsilon_3\varepsilon_4) & 1 - 2\varepsilon_3^2 - 2\varepsilon_1^2 & 2(\varepsilon_2\varepsilon_3 - \varepsilon_1\varepsilon_4) \\[4pt]
2(\varepsilon_3\varepsilon_1 - \varepsilon_2\varepsilon_4) & 2(\varepsilon_3\varepsilon_2 + \varepsilon_1\varepsilon_4) & 1 - 2\varepsilon_1^2 - 2\varepsilon_2^2
\end{bmatrix}
=
\begin{bmatrix}
a_{11} & a_{12} & a_{13} \\ a_{21} & a_{22} & a_{23} \\ a_{31} & a_{32} & a_{33}
\end{bmatrix}
\end{equation}

计算反对称矩阵：
\begin{equation}
\tilde{\boldsymbol{\omega}}_{ij|j} = A_{ij}^{\mathrm{T}}\dot{A}_{ij}
=
\begin{bmatrix}
a_{11} & a_{21} & a_{31} \\ a_{12} & a_{22} & a_{32} \\ a_{13} & a_{23} & a_{33}
\end{bmatrix}
\begin{bmatrix}
\dot{a}_{11} & \dot{a}_{12} & \dot{a}_{13} \\ \dot{a}_{21} & \dot{a}_{22} & \dot{a}_{23} \\ \dot{a}_{31} & \dot{a}_{32} & \dot{a}_{33}
\end{bmatrix}
=
\begin{bmatrix}
0                & -(\omega_z)_{ij|j} & (\omega_y)_{ij|j}   \\[2pt]
(\omega_z)_{ij|j} & 0                 & -(\omega_x)_{ij|j}  \\[2pt]
-(\omega_y)_{ij|j}& (\omega_x)_{ij|j} & 0
\end{bmatrix}
\end{equation}

利用约束方程及其导数：
\begin{align}
1 &= \varepsilon_1^2 + \varepsilon_2^2 + \varepsilon_3^2 + \varepsilon_4^2 \\
0 &= 2(\varepsilon_1\dot{\varepsilon}_1 + \varepsilon_2\dot{\varepsilon}_2 + \varepsilon_3\dot{\varepsilon}_3 + \varepsilon_4\dot{\varepsilon}_4)
\end{align}

\subsection{由角速度列阵反求欧拉参数对时间的一阶导数}

由式~(\ref{eq:7.8})，$\boldsymbol{\omega}_{ij|j}$的表达式可改写为增广形式：
\begin{equation}
\boldsymbol{\omega}_{ij|j} = 2
\begin{bmatrix}
\varepsilon_4  & \varepsilon_3  & -\varepsilon_2 & -\varepsilon_1 \\[2pt]
-\varepsilon_3 & \varepsilon_4  & \varepsilon_1  & -\varepsilon_2 \\[2pt]
\varepsilon_2  & -\varepsilon_1 & \varepsilon_4  & -\varepsilon_3
\end{bmatrix}
\begin{bmatrix}
\dot{\varepsilon}_1 \\ \dot{\varepsilon}_2 \\ \dot{\varepsilon}_3 \\ \dot{\varepsilon}_4
\end{bmatrix},
\qquad
0 = 2(\varepsilon_1\dot{\varepsilon}_1 + \varepsilon_2\dot{\varepsilon}_2 + \varepsilon_3\dot{\varepsilon}_3 + \varepsilon_4\dot{\varepsilon}_4)
\end{equation}

由此可得增广矩阵形式：
\begin{equation}
\begin{bmatrix}
\boldsymbol{\omega}_{ij|j} \\ 0
\end{bmatrix}
= 2
\begin{bmatrix}
\varepsilon_4  & \varepsilon_3  & -\varepsilon_2 & -\varepsilon_1 \\
-\varepsilon_3 & \varepsilon_4  & \varepsilon_1  & -\varepsilon_2 \\
\varepsilon_2  & -\varepsilon_1 & \varepsilon_4  & -\varepsilon_3 \\
\varepsilon_1  & \varepsilon_2  & \varepsilon_3  & \varepsilon_4
\end{bmatrix}
\begin{bmatrix}
\dot{\varepsilon}_1 \\ \dot{\varepsilon}_2 \\ \dot{\varepsilon}_3 \\ \dot{\varepsilon}_4
\end{bmatrix}
=: 2\bar{H}_{\mathrm{EP}}\dot{\boldsymbol{\varepsilon}}
\label{eq:7.10}
\end{equation}

其中
\begin{equation}
\bar{H}_{\mathrm{EP}} =
\begin{bmatrix}
\varepsilon_4  & \varepsilon_3  & -\varepsilon_2 & -\varepsilon_1 \\
-\varepsilon_3 & \varepsilon_4  & \varepsilon_1  & -\varepsilon_2 \\
\varepsilon_2  & -\varepsilon_1 & \varepsilon_4  & -\varepsilon_3 \\
\varepsilon_1  & \varepsilon_2  & \varepsilon_3  & \varepsilon_4
\end{bmatrix}
\end{equation}

可以验证$\bar{H}_{\mathrm{EP}}$为正交矩阵：
\begin{equation}
\bar{H}_{\mathrm{EP}}^{\mathrm{T}}\bar{H}_{\mathrm{EP}} =
\begin{bmatrix}
\varepsilon_4  & -\varepsilon_3 & \varepsilon_2  & \varepsilon_1  \\
\varepsilon_3  & \varepsilon_4  & -\varepsilon_1 & \varepsilon_2  \\
-\varepsilon_2 & \varepsilon_1  & \varepsilon_4  & \varepsilon_3  \\
-\varepsilon_1 & -\varepsilon_2 & -\varepsilon_3 & \varepsilon_4
\end{bmatrix}
\begin{bmatrix}
\varepsilon_4  & \varepsilon_3  & -\varepsilon_2 & -\varepsilon_1 \\
-\varepsilon_3 & \varepsilon_4  & \varepsilon_1  & -\varepsilon_2 \\
\varepsilon_2  & -\varepsilon_1 & \varepsilon_4  & -\varepsilon_3 \\
\varepsilon_1  & \varepsilon_2  & \varepsilon_3  & \varepsilon_4
\end{bmatrix}
= I_4
\end{equation}

因此，由角速度反求欧拉参数导数的表达式为
\begin{equation}
\dot{\boldsymbol{\varepsilon}} = \frac{1}{2}\bar{H}_{\mathrm{EP}}^{\mathrm{T}}
\begin{bmatrix}
\boldsymbol{\omega}_{ij|j} \\ 0
\end{bmatrix}
= \frac{1}{2}
\begin{bmatrix}
\varepsilon_4  & -\varepsilon_3 & \varepsilon_2  & \varepsilon_1  \\
\varepsilon_3  & \varepsilon_4  & -\varepsilon_1 & \varepsilon_2  \\
-\varepsilon_2 & \varepsilon_1  & \varepsilon_4  & \varepsilon_3  \\
-\varepsilon_1 & -\varepsilon_2 & -\varepsilon_3 & \varepsilon_4
\end{bmatrix}
\begin{bmatrix}
\boldsymbol{\omega}_{ij|j} \\ 0
\end{bmatrix}
= \frac{1}{2}
\begin{bmatrix}
\varepsilon_4 I_3 + \tilde{\bar{\boldsymbol{\varepsilon}}} \\[2pt]
-\bar{\boldsymbol{\varepsilon}}^{\mathrm{T}}
\end{bmatrix}
\boldsymbol{\omega}_{ij|j}
\label{eq:7.11}
\end{equation}

\section{欧拉角}

\subsection{体内ZXZ转动（经典欧拉角）}

从$i$系到$j$系的相对转动解读为绕连体坐标系（动系）坐标轴的三次相继转动，转轴分别为$Z$、$X$、$Z$，转角分别为$\theta_1$、$\theta_2$、$\theta_3$。在经典力学中：
\begin{itemize}
\item $\theta_1$：转轴$z_i$，进动角（$\psi$）
\item $\theta_2$：转轴$L$（节线），章动角（$\theta$）
\item $\theta_3$：转轴$z_j$，自转角（$\phi$）
\end{itemize}

坐标变换矩阵：
\begin{equation}
A_{ij} = A_z(\theta_1)A_x(\theta_2)A_z(\theta_3)
\end{equation}

展开得
\begin{align}
A_{ij} &= A_z(\theta_1)A_x(\theta_2)A_z(\theta_3) \nonumber\\
&= \begin{bmatrix}
\cos\theta_1 & -\sin\theta_1 & 0 \\
\sin\theta_1 & \cos\theta_1  & 0 \\
0            & 0            & 1
\end{bmatrix}
\begin{bmatrix}
1 & 0             & 0             \\
0 & \cos\theta_2  & -\sin\theta_2  \\
0 & \sin\theta_2  & \cos\theta_2
\end{bmatrix}
\begin{bmatrix}
\cos\theta_3  & -\sin\theta_3 & 0 \\
\sin\theta_3  & \cos\theta_3  & 0 \\
0            & 0            & 1
\end{bmatrix} \nonumber\\
&= \begin{bmatrix}
c_1c_3 - s_1c_2s_3  & -c_1s_3 - s_1c_2c_3  & s_1s_2    \\[3pt]
s_1c_3 + c_1c_2s_3  & c_1c_2c_3 - s_1s_3   & -c_1s_2   \\[3pt]
s_2s_3              & s_2c_3               & c_2
\end{bmatrix}
\label{eq:7.12}
\end{align}
其中$c_i = \cos\theta_i$，$s_i = \sin\theta_i$。

\subsubsection*{由坐标变换矩阵反求转角}

令
\begin{equation}
A_{ij} =
\begin{bmatrix}
c_1c_3 - s_1c_2s_3  & -c_1s_3 - s_1c_2c_3  & s_1s_2  \\
s_1c_3 + c_1c_2s_3  & c_1c_2c_3 - s_1s_3   & -c_1s_2 \\
s_2s_3              & s_2c_3               & c_2
\end{bmatrix}
=
\begin{bmatrix}
a_{11} & a_{12} & a_{13} \\ a_{21} & a_{22} & a_{23} \\ a_{31} & a_{32} & a_{33}
\end{bmatrix}
\end{equation}

\begin{enumerate}
\item[(I)] 当$c_2 = a_{33} \approx 1$时，$s_2 \approx 0$，此时
\begin{equation}
A_{ij} =
\begin{bmatrix}
c_1c_3 - s_1s_3 & -c_1s_3 - s_1c_3 & 0 \\[3pt]
s_1c_3 + c_1s_3 & c_1c_3 - s_1s_3  & 0 \\[3pt]
0               & 0                & 1
\end{bmatrix}
=
\begin{bmatrix}
a_{11} & a_{12} & 0 \\ a_{21} & a_{22} & 0 \\ 0 & 0 & 1
\end{bmatrix}
\end{equation}
其中$a_{11} = \cos(\theta_1 + \theta_3)$，$a_{21} = \sin(\theta_1 + \theta_3)$。可取
\begin{equation}
\theta_2 = 0 + 2k_1\pi, \quad
\theta_1 + \theta_3 = \operatorname{atan2}(a_{21}, a_{11}) + 2k_2\pi
\end{equation}

\item[(II)] 当$c_2 = a_{33} \approx -1$时，$s_2 \approx 0$，此时
\begin{equation}
A_{ij} =
\begin{bmatrix}
c_1c_3 + s_1s_3  & -c_1s_3 + s_1c_3 & 0  \\[3pt]
s_1c_3 - c_1s_3  & -c_1c_3 - s_1s_3 & 0  \\[3pt]
0                & 0                 & -1
\end{bmatrix}
=
\begin{bmatrix}
a_{11} & a_{12} & 0 \\ a_{21} & a_{22} & 0 \\ 0 & 0 & -1
\end{bmatrix}
\end{equation}
其中$a_{11} = \cos(\theta_1 - \theta_3)$，$a_{21} = \sin(\theta_1 - \theta_3)$。可取
\begin{equation}
\theta_2 = \pi + 2k_1\pi, \quad
\theta_1 - \theta_3 = \operatorname{atan2}(a_{21}, a_{11}) + 2k_2\pi
\end{equation}

\item[(III)] 当$|c_2| = |a_{33}| \neq 1$时，
\begin{equation}
\begin{bmatrix}
1   & -a_{33} \\ a_{33} & -1
\end{bmatrix}
\begin{bmatrix}
c_1c_3 \\ s_1s_3
\end{bmatrix}
=
\begin{bmatrix}
a_{11} \\ a_{22}
\end{bmatrix}
\quad\Rightarrow\quad
\begin{cases}
c_1c_3 = (a_{11} - a_{22}a_{33})/(1 - a_{33}^2) =: b_1 \\[3pt]
s_1s_3 = (-a_{22} + a_{11}a_{33})/(1 - a_{33}^2) =: b_2
\end{cases}
\end{equation}
\begin{equation}
\begin{bmatrix}
-a_{33} & -1 \\ -1 & a_{33}
\end{bmatrix}
\begin{bmatrix}
s_1c_3 \\ c_1s_3
\end{bmatrix}
=
\begin{bmatrix}
a_{12} \\ a_{21}
\end{bmatrix}
\quad\Rightarrow\quad
\begin{cases}
s_1c_3 = (a_{21} + a_{12}a_{33})/(1 - a_{33}^2) =: b_3 \\[3pt]
c_1s_3 = (-a_{12} - a_{21}a_{33})/(1 - a_{33}^2) =: b_4
\end{cases}
\end{equation}
进而有
\begin{equation}
\begin{cases}
\theta_1 + \theta_3 = \operatorname{atan2}(b_3 + b_4, b_1 - b_2) =: b_5 + 2k_1\pi \\[4pt]
\theta_1 - \theta_3 = \operatorname{atan2}(b_3 - b_4, b_1 + b_2) =: b_6 + 2k_2\pi
\end{cases}
\end{equation}
\begin{equation}
\begin{cases}
\theta_1 = (b_5 + b_6)/2 + k_3\pi \\[4pt]
\theta_3 = (b_5 - b_6)/2 - k_3\pi + 2k_4\pi
\end{cases}
\end{equation}
\end{enumerate}

\subsubsection*{体内ZXZ转动与角速度投影列阵的关系}

由$\tilde{\boldsymbol{\omega}}_{ij|j} = A_{ij}^{\mathrm{T}}\dot{A}_{ij}$可得：
\begin{equation}
\boldsymbol{\omega}_{ij|j} =
\begin{bmatrix}
\sin\theta_2\,\sin\theta_3 & \cos\theta_3  & 0 \\[3pt]
\sin\theta_2\,\cos\theta_3 & -\sin\theta_3 & 0 \\[3pt]
\cos\theta_2               & 0            & 1
\end{bmatrix}
\begin{bmatrix}
\dot{\theta}_1 \\ \dot{\theta}_2 \\ \dot{\theta}_3
\end{bmatrix}
=: H_{\mathrm{BZXZ}}\,\dot{\boldsymbol{\theta}}
\label{eq:7.13}
\end{equation}

系数矩阵的行列式为
\begin{equation}
\det(H_{\mathrm{BZXZ}}) = -\sin\theta_2
\end{equation}

因此，当$\theta_2 = k\pi$时，不存在从$\boldsymbol{\omega}_{ij|j}$到$\dot{\boldsymbol{\theta}}$的对应关系；也不存在从$\dot{\boldsymbol{\omega}}_{ij|j}$到$\ddot{\boldsymbol{\theta}}$的对应关系。$\theta_2 = k\pi$为此类三轴角的奇异点。

角速度列阵$\boldsymbol{\omega}_{ij|j}$对体内ZXZ转角$\boldsymbol{\theta}$的偏导数满足
\begin{equation}
\frac{\partial\boldsymbol{\omega}_{ij|j}}{\partial\boldsymbol{\theta}}
=
\begin{bmatrix}
0 & \cos\theta_2\sin\theta_3\,\dot{\theta}_1 & \sin\theta_2\cos\theta_3\,\dot{\theta}_1 - \cos\theta_3\,\dot{\theta}_2 \\[3pt]
0 & \cos\theta_2\cos\theta_3\,\dot{\theta}_1 & -\sin\theta_2\sin\theta_3\,\dot{\theta}_1 - \cos\theta_3\,\dot{\theta}_2 \\[3pt]
0 & -\sin\theta_2\,\dot{\theta}_1           & 0
\end{bmatrix}
\label{eq:7.14}
\end{equation}

\subsection{体内ZYX转动（卡尔丹角，Yaw-Pitch-Roll）}

从$i$系到$j$系的相对转动解读为绕连体坐标系（动系）坐标轴的三次相继转动，转轴分别为$Z$、$Y$、$X$，转角分别为$\theta_1$、$\theta_2$、$\theta_3$。在飞行器动力学中：
\begin{itemize}
\item $\theta_1$：转轴$z_i$，偏航角（$\psi$）
\item $\theta_2$：转轴$L$（节线），俯仰角（$\theta$）
\item $\theta_3$：转轴$x_j$，滚转角（$\phi$）
\end{itemize}

坐标变换矩阵：
\begin{equation}
A_{ij} = A_z(\theta_1)A_y(\theta_2)A_x(\theta_3)
\end{equation}

展开得
\begin{align}
A_{ij} &= A_z(\theta_1)A_y(\theta_2)A_x(\theta_3) \nonumber\\
&= \begin{bmatrix}
\cos\theta_1 & -\sin\theta_1 & 0 \\
\sin\theta_1 & \cos\theta_1  & 0 \\
0            & 0            & 1
\end{bmatrix}
\begin{bmatrix}
\cos\theta_2  & 0 & \sin\theta_2  \\
0            & 1 & 0             \\
-\sin\theta_2 & 0 & \cos\theta_2
\end{bmatrix}
\begin{bmatrix}
1 & 0            & 0             \\
0 & \cos\theta_3 & -\sin\theta_3  \\
0 & \sin\theta_3 & \cos\theta_3
\end{bmatrix} \nonumber\\
&= \begin{bmatrix}
c_1c_2  & -s_1c_3 + c_1s_2s_3  & s_1s_3 + c_1s_2c_3  \\[3pt]
s_1c_2  & c_1c_3 + s_1s_2s_3   & -c_1s_3 + s_1s_2c_3 \\[3pt]
-s_2    & c_2s_3               & c_2c_3
\end{bmatrix}
\label{eq:7.15}
\end{align}
其中$c_i = \cos\theta_i$，$s_i = \sin\theta_i$。

\subsubsection*{由坐标变换矩阵反求转角}

令
\begin{equation}
A_{ij} =
\begin{bmatrix}
c_1c_2  & -s_1c_3 + c_1s_2s_3  & s_1s_3 + c_1s_2c_3 \\
s_1c_2  & c_1c_3 + s_1s_2s_3   & -c_1s_3 + s_1s_2c_3 \\
-s_2    & c_2s_3               & c_2c_3
\end{bmatrix}
=
\begin{bmatrix}
a_{11} & a_{12} & a_{13} \\ a_{21} & a_{22} & a_{23} \\ a_{31} & a_{32} & a_{33}
\end{bmatrix}
\end{equation}

\begin{enumerate}
\item[(I)] 当$s_2 = -a_{31} \approx 1$时，$c_2 \approx 0$，此时
\begin{equation}
A_{ij} =
\begin{bmatrix}
0 & -s_1c_3 + c_1s_3 & s_1s_3 + c_1c_3 \\[3pt]
0 & c_1c_3 + s_1s_3  & -c_1s_3 + s_1c_3 \\[3pt]
-1 & 0               & 0
\end{bmatrix}
=
\begin{bmatrix}
0 & a_{12} & a_{13} \\ 0 & a_{22} & a_{23} \\ -1 & 0 & 0
\end{bmatrix}
\end{equation}
其中$a_{23} = \sin(\theta_1 - \theta_3)$，$a_{22} = \cos(\theta_1 - \theta_3)$。可取
\begin{equation}
\theta_2 = \pi/2 + 2k_1\pi, \quad
\theta_1 - \theta_3 = \operatorname{atan2}(a_{23}, a_{22}) + 2k_2\pi
\end{equation}

\item[(II)] 当$s_2 = -a_{31} \approx -1$时，$c_2 \approx 0$，此时
\begin{equation}
A_{ij} =
\begin{bmatrix}
0 & -s_1c_3 - c_1s_3 & s_1s_3 - c_1c_3 \\[3pt]
0 & c_1c_3 - s_1s_3  & -c_1s_3 - s_1c_3 \\[3pt]
1 & 0                & 0
\end{bmatrix}
=
\begin{bmatrix}
0 & a_{12} & a_{13} \\ 0 & a_{22} & a_{23} \\ 1 & 0 & 0
\end{bmatrix}
\end{equation}
其中$a_{23} = -\sin(\theta_1 + \theta_3)$，$a_{22} = \cos(\theta_1 + \theta_3)$。可取
\begin{equation}
\theta_2 = -\pi/2 + 2k_1\pi, \quad
\theta_1 + \theta_3 = \operatorname{atan2}(-a_{23}, a_{22}) + 2k_2\pi
\end{equation}

\item[(III)] 当$|s_2| = |a_{31}| \neq 1$时，
\begin{equation}
\begin{bmatrix}
1 & -a_{31} \\ -a_{31} & 1
\end{bmatrix}
\begin{bmatrix}
c_1c_3 \\ s_1s_3
\end{bmatrix}
=
\begin{bmatrix}
a_{22} \\ a_{13}
\end{bmatrix}
\quad\Rightarrow\quad
\begin{cases}
c_1c_3 = (a_{22} + a_{13}a_{31})/(1 - a_{31}^2) =: b_1 \\[3pt]
s_1s_3 = (a_{13} + a_{22}a_{31})/(1 - a_{31}^2) =: b_2
\end{cases}
\end{equation}
\begin{equation}
\begin{bmatrix}
-a_{31} & -1 \\ -1 & -a_{31}
\end{bmatrix}
\begin{bmatrix}
s_1c_3 \\ c_1s_3
\end{bmatrix}
=
\begin{bmatrix}
a_{23} \\ a_{12}
\end{bmatrix}
\quad\Rightarrow\quad
\begin{cases}
s_1c_3 = (-a_{12} + a_{23}a_{31})/(1 - a_{31}^2) =: b_3 \\[3pt]
c_1s_3 = (-a_{23} + a_{12}a_{31})/(1 - a_{31}^2) =: b_4
\end{cases}
\end{equation}
进而有
\begin{equation}
\begin{cases}
\theta_1 + \theta_3 = \operatorname{atan2}(b_3 + b_4, b_1 - b_2) =: b_5 + 2k_1\pi \\[4pt]
\theta_1 - \theta_3 = \operatorname{atan2}(b_3 - b_4, b_1 + b_2) =: b_6 + 2k_2\pi
\end{cases}
\end{equation}
\begin{equation}
\begin{cases}
\theta_1 = (b_5 + b_6)/2 + k_3\pi \\[4pt]
\theta_3 = (b_5 - b_6)/2 - k_3\pi + 2k_4\pi
\end{cases}
\end{equation}
\end{enumerate}

\subsubsection*{体内ZYX转动与角速度投影列阵的关系}

由$\tilde{\boldsymbol{\omega}}_{ij|j} = A_{ij}^{\mathrm{T}}\dot{A}_{ij}$可得：
\begin{equation}
\boldsymbol{\omega}_{ij|j} =
\begin{bmatrix}
-\sin\theta_2            & 0           & 1 \\[3pt]
\cos\theta_2\sin\theta_3 & \cos\theta_3 & 0 \\[3pt]
\cos\theta_2\cos\theta_3 & -\sin\theta_3& 0
\end{bmatrix}
\begin{bmatrix}
\dot{\theta}_1 \\ \dot{\theta}_2 \\ \dot{\theta}_3
\end{bmatrix}
=: H_{\mathrm{BZYX}}\,\dot{\boldsymbol{\theta}}
\label{eq:7.16}
\end{equation}

系数矩阵的行列式为
\begin{equation}
\det(H_{\mathrm{BZYX}}) = -\cos\theta_2
\end{equation}

因此，当$\theta_2 = \pi/2 + k\pi$时，不存在从$\boldsymbol{\omega}_{ij|j}$到$\dot{\boldsymbol{\theta}}$的对应关系；也不存在从$\dot{\boldsymbol{\omega}}_{ij|j}$到$\ddot{\boldsymbol{\theta}}$的对应关系。$\theta_2 = \pi/2 + k\pi$为此类三轴角的奇异点。

角速度列阵$\boldsymbol{\omega}_{ij|j}$对体内ZYX转角$\boldsymbol{\theta}$的偏导数满足
\begin{equation}
\frac{\partial\boldsymbol{\omega}_{ij|j}}{\partial\boldsymbol{\theta}}
=
\begin{bmatrix}
0 & -\cos\theta_2\,\dot{\theta}_1  & 0 \\[3pt]
0 & -\sin\theta_2\sin\theta_3\,\dot{\theta}_1 & \cos\theta_2\cos\theta_3\,\dot{\theta}_1 - \sin\theta_3\,\dot{\theta}_2 \\[3pt]
0 & -\sin\theta_2\cos\theta_3\,\dot{\theta}_1 & -\cos\theta_2\sin\theta_3\,\dot{\theta}_1 - \cos\theta_3\,\dot{\theta}_2
\end{bmatrix}
\label{eq:7.17}
\end{equation}

\section{坐标系姿态广义坐标的其他形式}

\begin{enumerate}
\item \textbf{旋转向量/罗德里格斯向量}
\begin{equation}
\boldsymbol{\rho} = \theta\boldsymbol{n}
\end{equation}

\item \textbf{罗德里格斯参数/吉布斯向量}（Rodrigues parameters / Gibbs vector）
\begin{equation}
\boldsymbol{g} = \boldsymbol{n}\tan\frac{\theta}{2}
\end{equation}

\item \textbf{修正的罗德里格斯参数}（Modified Rodrigues Parameters, MRP）
\begin{equation}
\boldsymbol{\sigma} = \boldsymbol{n}\tan\frac{\theta}{4}
\end{equation}
\end{enumerate}

\section{相继转动}

方向余弦矩阵与转轴转角的关系用转轴和转角表示为
\begin{equation}
A_{ij} = I_3\cos\theta + \tilde{\boldsymbol{n}}\sin\theta + \boldsymbol{n}\boldsymbol{n}^{\mathrm{T}}(1 - \cos\theta)
\end{equation}

当转轴$\vec{n}$相对于$i$系和$j$系固定时，
\begin{equation}
\dot{A}_{ij} = \bigl(-I_3\sin\theta + \tilde{\boldsymbol{n}}\cos\theta + \boldsymbol{n}\boldsymbol{n}^{\mathrm{T}}\sin\theta\bigr)\dot{\theta}
\end{equation}

接下来对$\tilde{\boldsymbol{\omega}}_{ij|j} = A_{ij}^{\mathrm{T}}\dot{A}_{ij}$进行化简：
\begin{align}
\tilde{\boldsymbol{\omega}}_{ij|j} &= A_{ij}^{\mathrm{T}}\dot{A}_{ij} \nonumber\\
&= \bigl[I_3\cos\theta + \tilde{\boldsymbol{n}}^{\mathrm{T}}\sin\theta + \boldsymbol{n}\boldsymbol{n}^{\mathrm{T}}(1 - \cos\theta)\bigr]
   \bigl(-I_3\sin\theta + \tilde{\boldsymbol{n}}\cos\theta + \boldsymbol{n}\boldsymbol{n}^{\mathrm{T}}\sin\theta\bigr)\dot{\theta} \nonumber\\
&= \bigl(-I_3\sin\theta\cos\theta + \tilde{\boldsymbol{n}}\cos^2\theta + \boldsymbol{n}\boldsymbol{n}^{\mathrm{T}}\sin\theta\cos\theta\bigr)\dot{\theta} \nonumber\\
&\quad + \bigl(-\tilde{\boldsymbol{n}}^{\mathrm{T}}\sin^2\theta + \tilde{\boldsymbol{n}}^{\mathrm{T}}\tilde{\boldsymbol{n}}\sin\theta\cos\theta + \cancel{\tilde{\boldsymbol{n}}^{\mathrm{T}}\boldsymbol{n}\boldsymbol{n}^{\mathrm{T}}\sin^2\theta}\bigr)\dot{\theta} \nonumber\\
&\quad + \bigl[-\boldsymbol{n}\boldsymbol{n}^{\mathrm{T}}\sin\theta(1 - \cos\theta)
       + \cancel{\boldsymbol{n}\boldsymbol{n}^{\mathrm{T}}\tilde{\boldsymbol{n}}\cos\theta(1 - \cos\theta)}
       + \boldsymbol{n}\boldsymbol{n}^{\mathrm{T}}\boldsymbol{n}\boldsymbol{n}^{\mathrm{T}}\sin\theta(1 - \cos\theta)\bigr]\dot{\theta} \nonumber\\
&= \tilde{\boldsymbol{n}}(\cos^2\theta + \sin^2\theta)\dot{\theta}
   + \bigl[-I_3 + \boldsymbol{n}\boldsymbol{n}^{\mathrm{T}} - (\boldsymbol{n}\boldsymbol{n}^{\mathrm{T}} - I_3)\bigr]\sin\theta\cos\theta\,\dot{\theta} \nonumber\\
&= \tilde{\boldsymbol{n}}\,\dot{\theta}
   = \widetilde{\boldsymbol{n}\dot{\theta}}
\end{align}

\begin{myformula}
由此可知，当转轴$\vec{n}$相对于$i$系和$j$系固定时，
\begin{equation}
\boldsymbol{\omega}_{ij|j} = \boldsymbol{n}\dot{\theta}
\label{eq:7.18}
\end{equation}
且
\begin{equation}
\boldsymbol{\omega}_{ij|i} = A_{ij}\boldsymbol{\omega}_{ij|j}
= \bigl[I_3\cos\theta + \tilde{\boldsymbol{n}}\sin\theta + \boldsymbol{n}\boldsymbol{n}^{\mathrm{T}}(1 - \cos\theta)\bigr]\boldsymbol{n}\dot{\theta}
= \boldsymbol{n}\dot{\theta} = \boldsymbol{\omega}_{ij|j}
\label{eq:7.19}
\end{equation}
\end{myformula}

当涉及多次相对转动时，为避免歧义，以上两式又可记为
\begin{equation}
\boldsymbol{\omega}_{ij|i} = \boldsymbol{\omega}_{ij|j} = \boldsymbol{n}_{ij}\dot{\theta}_{ij}
\end{equation}

设$l$系由$i$系经三次相继转动获得，相应的方向余弦矩阵表示为
\begin{equation}
A_{il} = A_{ij}(\boldsymbol{n}_{ij}, \theta_{ij})\,A_{jk}(\boldsymbol{n}_{jk}, \theta_{jk})\,A_{kl}(\boldsymbol{n}_{kl}, \theta_{kl})
\end{equation}
且$\boldsymbol{n}_{ij}$、$\boldsymbol{n}_{jk}$、$\boldsymbol{n}_{kl}$不随时间变化。此时，上式对时间的一阶导数可表示为
\begin{equation}
A_{il}\tilde{\boldsymbol{\omega}}_{il|l} = A_{ij}\tilde{\boldsymbol{\omega}}_{ij|j}A_{jk}A_{kl} + A_{ij}A_{jk}\tilde{\boldsymbol{\omega}}_{jk|k}A_{kl} + A_{ij}A_{jk}A_{kl}\tilde{\boldsymbol{\omega}}_{kl|l}
\end{equation}
式中
\begin{equation}
\boldsymbol{\omega}_{ij|j} = \boldsymbol{n}_{ij}\dot{\theta}_{ij}, \quad
\boldsymbol{\omega}_{jk|k} = \boldsymbol{n}_{jk}\dot{\theta}_{jk}, \quad
\boldsymbol{\omega}_{kl|l} = \boldsymbol{n}_{kl}\dot{\theta}_{kl}
\end{equation}

则$l$系相对于$i$系的角速度在$l$系中的投影为
\begin{equation}
\boldsymbol{\omega}_{il|l} = A_{jl}^{\mathrm{T}}\boldsymbol{\omega}_{ij|j} + A_{kl}^{\mathrm{T}}\boldsymbol{\omega}_{jk|k} + \boldsymbol{\omega}_{kl|l}
= A_{jl}^{\mathrm{T}}\boldsymbol{n}_{ij}\dot{\theta}_{ij} + A_{kl}^{\mathrm{T}}\boldsymbol{n}_{jk}\dot{\theta}_{jk} + \boldsymbol{n}_{kl}\dot{\theta}_{kl}
\label{eq:7.20}
\end{equation}
